\section{背景と目的}\label{sec:background}

表形式データの分類には、決定木や線形モデルなど複数の候補を同じ手順で比較する方法がある。AutoGluonのようにモデル選定やアンサンブルを自動化する基盤も使われている\cite{autogluon}。業務で使う際には、性能値の大小に加え、その値がどの対象、どの時点、どの分割で得られたかが判断に影響する。同じ顧客の記録が学習と評価にまたがるときと、新しい顧客を評価するときでは、同じF1でも答える問いが異なる。後続期間に適用する用途では、クラスの割合や入力分布が変わることもある。

本研究の目的は、分類器の選定を自動化し、その選定を支えたデータ点検と分割、単純ルールとの比較、誤分類の内訳をあとから検査できる評価基盤を構築することである。設計上の課題は三つある。第一に、記録の単位や予測時点を確認し、用途に合う分割を選べるようにする。第二に、評価値とその条件付き不確かさを出し、数値計算の差や測定・ラベルの偏りについて、どの範囲まで確認したかを明示する。第三に、単純ルールと同じ保留データで比べ、総合指標の差を見逃しと誤検知に分けて読めるようにする。

11月9日の短報は、モデル選定、固定モデルの条件付きブートストラップ、業務3例を中心に報告した。11月19日の改稿では、データ点検と抽出、グループ分割と日付順評価、Windowsの業務利用者向け画面を加えた\cite{letters,milestone}。本稿は両短報を踏まえ、各評価単位の設計、実験条件、結果の解釈、再実行の入口を一つにまとめる。主要な実証は教師あり分類で行い、バックエンド比較、特徴量探索、負荷測定を補足検証として扱う。
